\section{其他边界下的本征值}
\idx{本征值}{benzhengzhi}
两端固定的弦给出正弦族.换边界条件,分离常数的方程不变,变的是 $X$ 在端点上的条件,谱就换一族.下面都在 $0<x<\ell$ 上讨论 $X''+\lambda X=0$.三角函数族本身的正交和完备已经在傅里叶级数一章里证过,这里只列出分离变量要用的谱.

\subsection{可以写出来的谱}
第二类条件 $X'(0)=X'(\ell)=0$ 的非零解是余弦.$\lambda<0$ 时通解是指数,两端导数为零迫使两个系数都是零.$\lambda=0$ 时 $X=A+Bx$,$X'=B=0$,常数函数留下来.$\lambda=\omega^2>0$ 时 $X=\cos(\omega x)$ 满足左端,$X'(\ell)=0$ 要求 $\sin(\omega\ell)=0$,
\begin{equation}
  \lambda_n=\Bigl(\frac{n\pi}{\ell}\Bigr)^2,
  \qquad
  X_n=\cos\Bigl(\frac{n\pi x}{\ell}\Bigr),
  \qquad n=0,1,2,\ldots.
\end{equation}
$n=0$ 是整段一起平移的模式,热方程和波动方程里它不衰减,也不振动.

一端固定、一端自由,$X(0)=X'(\ell)=0$.$\lambda\leq 0$ 没有非零解.$\lambda=\omega^2$ 时 $X=\sin(\omega x)$,右端导数 $-\omega\cos(\omega\ell)=0$,所以 $\omega\ell=(n+\half)\pi$,
\begin{equation}
  \lambda_n=\Bigl(\frac{(n+\half)\pi}{\ell}\Bigr)^2,
  \qquad
  X_n=\sin\Bigl(\frac{(n+\half)\pi x}{\ell}\Bigr),
  \qquad n=0,1,2,\ldots.
  \label{eq:half-integer}
\end{equation}
两端对调,$X'(0)=X(\ell)=0$,同一组 $\lambda_n$,本征函数改成 $\cos\bigl((n+\half)\pi x/\ell\bigr)$.这就是半整数正弦族或余弦族.在 $[0,\ell]$ 上它们的范数仍是 $\ell/2$.

周期条件 $X(0)=X(\ell)$, $X'(0)=X'(\ell)$ 描述环或晶体原胞.$\lambda=0$ 只有常数.$\lambda>0$ 时两个独立解一起留下来,本征值二重,
\begin{equation}
  \lambda_n=\Bigl(\frac{2\pi n}{\ell}\Bigr)^2,
  \qquad
  X=\cos\Bigl(\frac{2\pi n x}{\ell}\Bigr),\;
  \sin\Bigl(\frac{2\pi n x}{\ell}\Bigr),
\end{equation}
$n=1,2,\ldots$.写成复指数 $e^{2\pi\imath n x/\ell}$,$n\in\mathbb{Z}$,每个 $n\neq 0$ 的模长相同,简并还在,只是基换了.

准周期条件把 Bloch 相位写进端点.设准动量 $\kappa\in(0,2\pi/\ell)$,
\begin{equation}
  X(\ell)=e^{\imath\kappa\ell}X(0),
  \qquad
  X'(\ell)=e^{\imath\kappa\ell}X'(0).
\end{equation}
解是 $X_n=\exp\bigl(\imath(2\pi n/\ell+\kappa)x\bigr)$,本征值 $(2\pi n/\ell+\kappa)^2$,$n\in\mathbb{Z}$,一般不再简并,本征函数也不再是实的.$\kappa=0$ 回到周期条件.$\kappa=\pi/\ell$ 是反周期,$X(\ell)=-X(0)$.自由电子在周期势里的能带,就是把这里的 $X''+\lambda X=0$ 换成 $X''+(\lambda-V(x))X=0$.

\begin{example}
左端固定、右端自由的热棒,初温 $g(x)=x$.写出温度.
\end{example}
\begin{solution}
用\eqref{eq:half-integer}.$u_t=k u_{xx}$,$u(0,t)=u_x(\ell,t)=0$.
\begin{equation}
  u(x,t)=\sum_{n=0}^{\infty}A_n
  \exp\Bigl(-k\lambda_n t\Bigr)
  \sin\Bigl(\frac{(n+\half)\pi x}{\ell}\Bigr).
\end{equation}
范数是 $\ell/2$,
\begin{equation}
  A_n=\frac{2}{\ell}\int_0^{\ell}x\sin\Bigl(\frac{(n+\half)\pi x}{\ell}\Bigr)\,\mathrm{d}x
  =\frac{2\ell(-1)^n}{(n+\half)^2\pi^2}.
\end{equation}
$t\to+\infty$ 时每一项都衰减,最低模式 $n=0$ 衰减最慢,速率是 $k\pi^2/(4\ell^2)$.
\end{solution}

\subsection{两端都是第三类}
\idx{罗宾条件}{luobin}
\begin{equation}
  X'(0)=\alpha X(0),
  \qquad
  X'(\ell)=-\beta X(\ell).
\end{equation}
$\alpha,\beta$ 同号时,弹性能的符号决定谱落在哪一边.分部积分,
\begin{equation}
  \lambda\int_0^{\ell}X^2\,\mathrm{d}x
  =\int_0^{\ell}(X')^2\,\mathrm{d}x+\alpha X(0)^2+\beta X(\ell)^2.
  \label{eq:robin-energy}
\end{equation}
因此 $\alpha\geq 0$, $\beta\geq 0$ 且不同时为零时,$\lambda>0$.令 $\lambda=\omega^2$,$X=\omega\cos(\omega x)+\alpha\sin(\omega x)$ 已经满足左端.代入右端,
\begin{equation}
  \tan(\omega\ell)=\frac{(\alpha+\beta)\omega}{\omega^2-\alpha\beta},
  \qquad
  X=\omega\cos(\omega x)+\alpha\sin(\omega x).
  \label{eq:robin-both}
\end{equation}
$\alpha,\beta\to 0$ 时根靠近 $(n-1)\pi/\ell$,回到余弦族;两者都趋于 $+\infty$ 时根靠近 $n\pi/\ell$,回到正弦族.

$\alpha$ 或 $\beta$ 为负时,弹性能可以是负的,\eqref{eq:robin-energy}不再禁止 $\lambda<0$.令 $\lambda=-\gamma^2$,$\gamma>0$,
\begin{equation}
  \tanh(\gamma\ell)=-\frac{(\alpha+\beta)\gamma}{\gamma^2+\alpha\beta},
  \qquad
  X=\gamma\cosh(\gamma x)+\alpha\sinh(\gamma x).
\end{equation}
$\lambda=0$ 是否是本征值,用直线 $X=cx+d$ 代进去:非零解存在当且仅当
\begin{equation}
  \alpha+\beta+\alpha\beta\ell=0.
  \label{eq:robin-threshold}
\end{equation}
这条双曲线把 $(\alpha,\beta)$ 平面分成若干块.把 $\alpha$ 或 $\beta$ 增大,\eqref{eq:robin-energy}的二次型只增不减,本征值因此不减.特别地 $\alpha=\beta$ 时可以把区间平移到中点,本征函数分成偶函数和奇函数.偶的一支在 $\alpha<0$ 时恰有一个负本征值,奇的一支还要求 $\alpha\ell<-2$.所以
\begin{itemize}
  \item $\alpha\geq 0$:没有负本征值;
  \item $-2<\alpha\ell<0$:恰有一个,本征函数为偶;
  \item $\alpha\ell<-2$:恰有两个,一偶一奇.
\end{itemize}
$\alpha\ell=-2$ 和 $\alpha=0$ 正好落在\eqref{eq:robin-threshold}上,零本征值出现在两块区域的边界.

\subsection{梁的横振动}
细梁的小挠度满足四阶方程
\begin{equation}
  u_{tt}+K u_{xxxx}=0, \qquad K>0.
\end{equation}
分离 $u=X(x)T(t)$ 得到 $X^{(4)}=\lambda X$, $T''+K\lambda T=0$.两端夹紧是 $X=X'=0$.能量
\begin{equation}
  \int_0^{\ell}(X'')^2\,\mathrm{d}x
\end{equation}
在这个边界下为正,所以 $\lambda=k^4>0$.通解
\begin{equation}
  X=A\cosh(kx)+B\sinh(kx)+C\cos(kx)+D\sin(kx).
\end{equation}
左端夹紧给出 $C=-A$, $D=-B$,
\begin{equation}
  X=A\bigl(\cosh(kx)-\cos(kx)\bigr)+B\bigl(\sinh(kx)-\sin(kx)\bigr).
\end{equation}
右端 $X(\ell)=X'(\ell)=0$ 是关于 $A,B$ 的齐次方程组,行列式为零化成
\begin{equation}
  \cosh(k\ell)\cos(k\ell)=1.
  \label{eq:beam-clamped}
\end{equation}
左端夹紧、右端自由时,自由端是弯矩和剪力为零,$X''(\ell)=X'''(\ell)=0$,同样的消元换成
\begin{equation}
  \cosh(k\ell)\cos(k\ell)=-1.
\end{equation}
两端都自由时谱和两端夹紧相同,但本征函数改成 $\cosh+\cos$ 与 $\sinh+\sin$ 的组合,并且 $\lambda=0$ 是二重本征值,刚体模式是 $1$ 和 $x$.

\begin{example}
说明两端夹紧时\eqref{eq:beam-clamped}的根都是正的,并且 $k\ell$ 落在 $(n\pi,(n+\half)\pi)$ 附近.
\end{example}
\begin{solution}
$\cosh(k\ell)\geq 1$,所以 $\cos(k\ell)=1/\cosh(k\ell)\in(0,1]$.$k=0$ 使 $X$ 恒为零,不是本征函数.在每个 $(2m\pi,(2m+1/2)\pi)$ 里余弦从 $1$ 降到 $0$,$1/\cosh$ 从略小于 $1$ 降到更小的正数,恰有一个交点.$(2m+3/2)\pi$ 到 $(2m+2)\pi$ 里余弦从 $0$ 升到 $1$,又有一个交点.大 $k$ 时 $1/\cosh$ 几乎是零,根靠近余弦的零点,也就是半整数个 $\pi$.
\end{solution}

\section{热方程的分离变量}
\idx{热方程}{refangcheng}
$u_t=k u_{xx}$,两端为零,$u(x,0)=g(x)$.分离后 $T'=-k\lambda T$,时间因子是衰减指数,不是振动,
\begin{equation}
  u(x,t)=\sum_{n=1}^{\infty}A_n e^{-k(n\pi/\ell)^2 t}
  \sin\Bigl(\frac{n\pi x}{\ell}\Bigr),
  \qquad
  A_n=\frac{2}{\ell}\int_0^{\ell}g(x)\sin\Bigl(\frac{n\pi x}{\ell}\Bigr)\,\mathrm{d}x.
  \label{eq:heat-series}
\end{equation}
$t<0$ 时指数变成增长,高模式放得最大,初值问题向过去不适定.$t\to+\infty$ 时
\begin{equation}
  u(x,t)=A_1 e^{-k\pi^2 t/\ell^2}\sin\Bigl(\frac{\pi x}{\ell}\Bigr)
  +O\bigl(e^{-4k\pi^2 t/\ell^2}\bigr).
\end{equation}
最低模式决定长时间的形状.

若方程和边界都不随时间变化,
\begin{equation}
  u_t=k u_{xx}+f(x),
  \qquad u(0,t)=\varphi, \quad u(\ell,t)=\psi,
\end{equation}
先丢掉初值,求稳态 $v''=-f/k$, $v(0)=\varphi$, $v(\ell)=\psi$.$u-v$ 满足齐次边界,按\eqref{eq:heat-series}衰减,所以
\begin{equation}
  u(x,t)=v(x)+O\bigl(e^{-k\pi^2 t/\ell^2}\bigr).
\end{equation}
温度趋于稳态.两端若改成齐次诺伊曼或周期条件,$\lambda_1=0$,常数模式不衰减.这时稳态存在的条件是总热流为零:两端诺伊曼 $u_x(0)=\varphi$, $u_x(\ell)=\psi$ 时
\begin{equation}
  \frac{1}{k}\int_0^{\ell}f\,\mathrm{d}x-\varphi+\psi=0,
\end{equation}
周期条件时积分本身必须是零.热流不平衡时,多出来的热量均匀堆积,$u=pt+A_1+O(e^{-k\lambda_2 t})$,$p$ 是 $f$ 的平均值.

\begin{example}
$u_t=u_{xx}$, $0<x<1$, $u(0,t)=0$, $u(1,t)=1$, $u(x,0)=0$.
\end{example}
\begin{solution}
稳态是 $v=x$.令 $w=u-x$,则 $w_t=w_{xx}$, $w(0,t)=w(1,t)=0$, $w(x,0)=-x$.
\begin{equation}
  A_n=2\int_0^1(-x)\sin(n\pi x)\,\mathrm{d}x=\frac{2(-1)^n}{n\pi},
\end{equation}
\begin{equation}
  u(x,t)=x+\sum_{n=1}^{\infty}\frac{2(-1)^n}{n\pi}e^{-n^2\pi^2 t}\sin(n\pi x).
\end{equation}
$t\to+\infty$ 时级数消失,温度变成线性分布 $x$.
\end{solution}

\section{薛定谔方程与零本征值的时间因子}
有限区间上的自由薛定谔方程 $u_t=\imath k u_{xx}$,边界仍是分离变量能接受的那几类.空间本征值照旧,$T'=-\imath k\lambda T$,
\begin{equation}
  u(x,t)=\sum_n A_n e^{-\imath k\lambda_n t}X_n(x).
\end{equation}
指数的模是 $1$.向过去和向未来都适定,但没有热方程那种趋于稳态的衰减.

波动方程 $u_{tt}=c^2 u_{xx}$ 在同一组 $X_n$ 上,
\begin{equation}
  T_n(t)=
  \begin{cases}
    A_n\cos(c\sqrt{\lambda_n}\,t)+B_n\sin(c\sqrt{\lambda_n}\,t), & \lambda_n>0,\\
    A_n+B_n t, & \lambda_n=0,\\
    A_n\cosh(c\sqrt{-\lambda_n}\,t)+B_n\sinh(c\sqrt{-\lambda_n}\,t), & \lambda_n<0.
  \end{cases}
\end{equation}
$\lambda=0$ 的线性项是刚体漂移,出现在两端自由或周期条件里.$\lambda<0$ 的双曲函数随时间增长,对应弹性支承吸能、平衡不稳定.系数仍由 $u(x,0)$ 和 $u_t(x,0)$ 的展开定,$B_n$ 在三种情形下分别除以 $c\sqrt{\lambda_n}$、$1$ 和 $c\sqrt{-\lambda_n}$.

\section{半带和矩形}
\idx{半带}{bandai}
半带 $0<x<\ell$, $y>0$ 上 $\Delta u=0$,左右是齐次边界,底边 $u(x,0)=g(x)$,并且 $u$ 有界.分离后 $X$ 仍是上面的本征函数,$Y''=\lambda Y$.$\lambda_n>0$ 时 $e^{\sqrt{\lambda_n} y}$ 无界,丢掉,
\begin{equation}
  u(x,y)=\sum_n A_n e^{-\sqrt{\lambda_n}\,y}X_n(x),
\end{equation}
$A_n$ 是 $g$ 的系数.若在 $y=0$ 再指定 $u_y$,两个指数都要留下,高频模式会被放大,柯西问题不适定.$\lambda=0$ 时 $Y=A+By$,有界解只能留常数.四周都是诺伊曼条件时,解存在还要求 $\int g X_0\,\mathrm{d}x=0$.

矩形 $0<x<a$, $0<y<b$,左右齐次,上下给定 $u(x,0)=g(x)$, $u(x,b)=h(x)$.两个指数都留在有限区间里,把它们收成双曲正弦后在端点取到 $1$,
\begin{equation}
  u(x,y)=\sum_n
  \Biggl(
    g_n\frac{\sinh\bigl(\sqrt{\lambda_n}(b-y)\bigr)}{\sinh(\sqrt{\lambda_n}b)}
    +h_n\frac{\sinh(\sqrt{\lambda_n}y)}{\sinh(\sqrt{\lambda_n}b)}
  \Biggr)X_n(x).
  \label{eq:rectangle-sinh}
\end{equation}
上下也改成诺伊曼、并且 $\lambda_0=0$ 时,解存在的条件是上下边界的总通量互相抵消.四边都是诺伊曼时,存在条件是四条边上的法向导数积分和为零,也就是没有净热流.函数 $v=x^2-y^2$ 本身调和,可以用来把两条边上的相容条件和另外两条拆开,再叠加回去.

\begin{example}
$0<x<\pi$, $0<y<b$, $u(0,y)=u(\pi,y)=u(x,0)=0$, $u(x,b)=\sin x$.
\end{example}
\begin{solution}
左右为零,所以 $X_n=\sin(nx)$, $\lambda_n=n^2$.顶边只有 $n=1$.\eqref{eq:rectangle-sinh}只剩一项,
\begin{equation}
  u(x,y)=\frac{\sinh y}{\sinh b}\sin x.
\end{equation}
底边 $\sinh 0=0$,顶边比值是 $1$.直接计算 $\Delta u=0$.
\end{solution}

\section{矩形域上的乘积解}
矩形膜 $0<x<a$, $0<y<b$,四周固定,$u_{tt}=c^2\Delta u$.令 $u=X(x)Y(y)T(t)$,三个变量各自成常数,
\begin{equation}
  X''+\lambda X=0, \quad Y''+\mu Y=0, \quad T''+c^2(\lambda+\mu)T=0.
\end{equation}
四周为零给出
\begin{equation}
  \lambda_m=\Bigl(\frac{m\pi}{a}\Bigr)^2, \quad
  \mu_n=\Bigl(\frac{n\pi}{b}\Bigr)^2, \quad
  \omega_{mn}=c\pi\sqrt{\frac{m^2}{a^2}+\frac{n^2}{b^2}},
\end{equation}
$m,n\geq 1$.一般解是
\begin{equation}
  u=\sum_{m,n}\bigl(A_{mn}\cos\omega_{mn}t+B_{mn}\sin\omega_{mn}t\bigr)
  \sin\Bigl(\frac{m\pi x}{a}\Bigr)\sin\Bigl(\frac{n\pi y}{b}\Bigr).
\end{equation}
热方程把 $\cos$ 和 $\sin$ 换成 $e^{-k(\lambda_m+\mu_n)t}$.长方体是同一件事再乘一个 $z$ 方向的正弦.正方形 $a=b$ 时 $\omega_{mn}=\omega_{nm}$,不同的 $(m,n)$ 可以共用一个频率.

稳态问题 $\Delta u=0$ 不能让 $\lambda$ 和 $\mu$ 同时为正又同时满足四边为零,否则 $T$ 的角色由一个会在某个方向上指数增长的因子扮演.四边为零的调和函数只有 $u=0$.非齐次边界才用上一节的双曲正弦.

\section{不先分离变量的展开}
有的方程对 $x$ 和 $t$ 分离不开,边界却仍然挑出正弦族.这时先把未知函数按这条边界允许的基展开,系数是 $t$ 的函数.

\begin{example}
$u_{tt}-2a u_{tx}-u_{xx}=0$, $0<x<\pi$, $u(0,t)=u(\pi,t)=0$.
\end{example}
\begin{solution}
试 $u=e^{\imath k(t-ax)}\sin(nx)$.代入后 $\sin(nx)$ 的系数是 $n^2-k^2(1+a^2)$.它为零时
\begin{equation}
  k=\pm\frac{n}{\sqrt{1+a^2}}, \qquad n=1,2,\ldots.
\end{equation}
两端的正弦自动满足边界.一般解是这两支的叠加,
\begin{equation}
  u=\sum_{n=1}^{\infty}
  \Bigl(
    A_n\cos\frac{n(t-ax)}{\sqrt{1+a^2}}
    +B_n\sin\frac{n(t-ax)}{\sqrt{1+a^2}}
  \Bigr)\sin(nx).
\end{equation}
$t=0$ 时 $e^{0}=1$,$A_n,B_n$ 仍由初值的正弦系数定.相因子 $t-ax$ 是输运项 $u_{tx}$ 留下的,普通的 $X(x)T(t)$ 乘积接不住它.
\end{solution}

矩形上的双调和方程 $\Delta^2 u=0$,若 $x=0,\pi$ 上 $u=u_{xx}=0$,同样可以先写 $u=\sin(nx)Y(y)$.$Y$ 满足 $(D^2-n^2)^2 Y=0$,通解是 $(A+By)e^{ny}+(C+Dy)e^{-ny}$,再由 $y$ 方向的两个边界条件定系数.

\section{多重傅里叶级数}
每个变量各自以 $2\pi$ 为周期时,系数可以一个变量一个变量地算.二元情形
\begin{equation}
  u(x,y)=\sum_{m,n\in\mathbb{Z}}c_{mn}e^{\imath(mx+ny)},
  \qquad
  c_{mn}=\frac{1}{(2\pi)^2}\int_0^{2\pi}\int_0^{2\pi}
  u(x,y)e^{-\imath(mx+ny)}\,\mathrm{d}x\,\mathrm{d}y.
\end{equation}
帕塞瓦尔等式是 $\sum|c_{mn}|^2=(2\pi)^{-2}\iint|u|^2$.格子 $\Gamma$ 换成一般的周期格时,先把基本胞腔线性变到立方体,对偶格 $\Gamma^*$ 上的波矢代替 $(m,n)$.分离变量在矩形上得到的正弦乘积,就是这个级数在齐次边界下的特例.
